\vspace{0.5em}\noindent\textbf{The problem}\par

An online store may carry thousands of products, while a typical shopper
only has the patience to browse a few dozen. Ranking products by sales
volume or upload date means most of the catalogue is never seen, and
users have to hunt for what they want on their own.

The system in this workshop addresses that by learning from actual user
behaviour to suggest products tailored to each person.

\vspace{0.5em}\noindent\textbf{Architecture}\par

\begin{figure}
\centering
\pandocbounded{\includegraphics[keepaspectratio,alt={Overall architecture}]{5-Workshop/architecture.jpg}}
\caption{Overall architecture}
\end{figure}

The system splits into two independent paths:

\textbf{The content delivery path.} Users hit the CloudFront domain, and
CloudFront serves the pre-built HTML, CSS and JavaScript files stored in
an S3 bucket. Because this is static content, it can be cached at edge
locations close to users, cutting latency and incurring almost no
compute cost.

\textbf{The data path.} The frontend calls API Gateway, which invokes a
Lambda function. Depending on the route, Lambda queries DynamoDB for
business data or calls Amazon Personalize for a recommendation list.

\vspace{0.5em}\noindent\textbf{Why these services}\par

\textbf{Serverless instead of traditional servers.} Traffic for an
internship project is low and uneven. With EC2 or ECS you pay by the
hour even when nobody visits. With Lambda, no requests means no cost.

\textbf{DynamoDB instead of RDS.} RDS bills per instance-hour even when
idle. DynamoDB in on-demand mode bills only for actual reads and writes.
For key-based data with simple access patterns, DynamoDB is the better
fit.

\textbf{Amazon Personalize instead of training on SageMaker.}
Implementing a recommendation algorithm from scratch requires deep
expertise and a lot of tuning time. Personalize packages the same
algorithms Amazon uses in its own products, letting you focus on data
quality rather than model implementation.

\textbf{CloudFront in front of S3.} Beyond speed, CloudFront lets you
keep the bucket fully private and provides HTTPS at no extra cost.

\vspace{0.5em}\noindent\textbf{Design principles applied
throughout}\par

\begin{itemize}
\tightlist
\item
  \textbf{Never trust data from the browser.} The server always looks up
  prices and recalculates totals itself, rather than using figures sent
  by the frontend.
\item
  \textbf{Least privilege.} Each component receives only the permissions
  it needs, on only the resources it needs.
\item
  \textbf{No hard-coded credentials.} Lambda uses an execution role; no
  access keys are embedded in source code.
\end{itemize}
